\chapter{在工程中使用RVM}\label{ch:6-using-rvm}

\section{总览}\label{sec:6-1}

本章节将描述在工程中使用RVM及其工程生成器的一般流程。该流程首先将功能划分到保护域中，然后编写工程描述文件以反应这一划分，最后对生成的工程做一定的调整以追加行为细节。

受限于篇幅，本章将仅描述部分重要配置项的设置；对于其他配置项，请参阅前面的各章节。

\section{软件组件到保护域的划分}\label{sec:6-2}

与裸机开发不同，在使用RVM的工程中，软件的各个组件将被指派到不同的保护域中，各个保护域的权限都受限以约束其内部软件的行为。一旦其中一个保护域因编码故障而崩溃，其余的保护域仍将保持运行，不会受到干扰；即便有保护域被黑客攻击劫持，黑客也只能获得该保护域的权限，无法控制整个系统。

决定如何划分软件中的各组件到合适的保护域中至关重要。至少要考虑以下三个方面：（1）组件划分的粒度，（2）组件对兼容性和实时性的要求和（3）组件对设备访问的要求。为了更好地把控系统的性能，推荐先在目标微控制器上运行各项基准测试，待测试完成后视各操作开销决定组件的划分粒度。

\subsection{组件划分的粒度}\label{subsec:6-2-1}

组件划分的粒度指独立组件的大小。理论上讲，任何软件都是无限可分的，直到划分为单条指令为止，而且划分粒度越细，其保护效果越好\footnote{极限划分的保护域约等于数据流安全}。然而，划分粒度越细，跨保护域通信就越多，而跨保护域通信是有代价的\footnote{这些代价可以在基准测试中看到}。因此，推荐将组件划分为合适的大小以尽量规避这一开销。合理的组件划分需要满足高内聚、低耦合这一原则，组件内部的联系应当紧密，而组件间的联系应当稀疏。换言之，一个组件应当负责一件不可或难以分割的功能，该功能的所有状态和逻辑都位于该组件之内，且该组件对其它组件提供的接口应当简明扼要、易于利用。

在划分开始时，可以将软件的每项独立功能都尽可能细地划分为组件，然后观察划分中是否有控制流和数据流的紧密耦合。有紧密耦合的组件具备如下特征之一：

1.
两个组件间互相调用的频次远高于其它组件的平均水平，比如在一个计算循环中反复调用另一组件；一旦两个组件被放在独立的保护域内，其通信开销将无法忍受。

2.
两个组件共享某个数据结构，且都有可能对该数据结构发起难以隔离、难以原子化的写操作，比如图变换操作、数据排序等；将这两个组件放在独立的保护域内没有意义，因为一个组件的崩溃必然招致数据结构的不完整，从而导致另一个组件崩溃。

3.
两个组件可以被视为一个更好描述的天然功能的拆分，合并后该组件的描述反而更加简洁；如果是这种情况，将这两个组件分配到独立的保护域中意义可能有限，因为它们不是一个逻辑功能，不需要具备崩溃后的独立恢复能力。

如果发现现有划分中存在这样的紧密耦合，两个组件即可合并为一个组件；当所有紧密耦合的组件都被合并后得到的就是一个合理的组件分割。在完成初始整合后若性能仍不达预期，则可以继续合并组件，直到得到安全性和性能之间的良好权衡。

\subsection{组件对兼容性和实时性的要求}\label{subsec:6-2-2}

在完成组件到保护域的划分后，需要进一步决定保护域的形式。RVM提供两种保护域，分别是原生进程和虚拟机，两者各有优劣。

在兼容性方面，原生进程必须使用RME原生系统调用进行编程，对其它运行时环境的兼容性则欠佳，如需使用则必须手动移植；虚拟机则可灵活使用任何客户机RTOS，并直接兼容该RTOS的一切生态环境。在实时性方面，原生进程无需虚拟机监视器作为中介，跨保护域通信性能更好，中断延迟更低，可以拥有与裸机RTOS相仿的性能；虚拟机则相反，跨保护域通信需要虚拟机监视器介导，中断延迟时间大约是裸机的3-4倍。

几个常见的例子如下表所述。

\begin{xltabular}{\textwidth}{@{}>{\hsize=1.0000\hsize\raggedright\arraybackslash}X>{\hsize=1.0000\hsize\raggedright\arraybackslash}X>{\hsize=1.0000\hsize\raggedright\arraybackslash}X>{\hsize=1.0000\hsize\raggedright\arraybackslash}X@{}}
\caption{组件划分样例}\label{tab:6-1}\\
\toprule
组件名称 & 组件划分 & 原因 & 代表性软件\tabularnewline
\endfirsthead
\toprule
组件名称 & 组件划分 & 原因 & 代表性软件\tabularnewline
\midrule
\endhead
\bottomrule
\endlastfoot
键控遥控码解析\strut
 & 原生进程\strut
 & 对实时性要求高，代码简单，需要准确的定时器嘀嗒。\strut
 & PT2262/PT2272解析

38kHz红外解析\strut
\tabularnewline
电力拖动驱动\strut
 & 原生进程\strut
 & 对实时性要求高，代码简单，且不依赖于RTOS。\strut
 & 步进/伺服电机驱动

变频器开关管控制驱动\strut
\tabularnewline
单线总线驱动\strut
 & 原生进程\strut
 & 对实时性要求高，代码简单，且不依赖于RTOS。\strut
 & DS18B20

ATECC608A\strut
\tabularnewline
音频总线驱动\strut
 & 原生进程\strut
 & 对实时性要求高，代码简单，需要及时刷新音频缓冲区。\strut
 & 各类I2S声卡驱动

各类SPDIF输出驱动\strut
\tabularnewline
梯形图执行器 & 原生进程 &
虽然代码复杂，但往往经过严格审查，且需要周期性地执行，必须移植到原生进程。
& OpenPLC\tabularnewline
网络协议栈\strut
 & 虚拟机\strut
 & 对实时性要求不高，但内部逻辑较复杂，需要RTOS支持。\strut
 & LwIP

uIP\strut
\tabularnewline
图形界面\strut
 & 虚拟机\strut
 & 对实时性要求不高，且可能依赖于RTOS支持。\strut
 & uC/GUI

LVGL\strut
\tabularnewline
文件系统\strut
 & 虚拟机\strut
 & 对实时性要求不高，且可能依赖于RTOS支持。\strut
 & ElmChan-FatFS

Yaffs\strut
\tabularnewline
脚本解释器\strut
 & 虚拟机\strut
 & 对实时性要求不高，且可能依赖于RTOS支持。\strut
 & Micropython

Lua\strut
\tabularnewline
串行存储器驱动\strut
 & 虚拟机\strut
 & 该驱动虽然简单，但对实时性要求不高，可以放在虚拟机中。\strut
 & AT24C02驱动

W25Q64驱动\strut
\tabularnewline
\end{xltabular}

\subsection{组件对设备访问的要求}\label{subsec:6-2-3}

部分驱动组件可能需要访问设备，因此需要将设备暴露给该组件所在的保护域。最常见也是最简单的做法是（1）仅将需要的设备设置为对该保护域可见，这可以通过将设备的内存映射加入该保护域的内存描述来实现。如果（2）多个保护域都需要使用同一个设备，则可将该设备驱动封装在一个单独的保护域中，其它保护域可以委托它来发起传输。

对于某些有较强自主性的设备，采用方法（1）将它们直接暴露给不受信任的保护域会导致重大信息安全隐患。虽然保护域内的代码没有能力突破时空保护域，但它可以配置这些自主设备，并将它们作为跳板来访问任意内存地址，或者生成高频次中断而发起拒绝服务攻击。最常见的此类设备是直接内存访问（Direct
Memory
Access，DMA）和定时器；恶意代码可以操纵前者来修改内核空间，或操纵后者来生成高频中断。

对于DMA等会自主访问内存的硬件，推荐的方法是将其交给一个单独的受信任的保护域，其它保护域需要使用DMA时必须委托受信保护域进行。这虽然增加了DMA的延迟，但保证无法利用DMA来破坏隔离边界。当然，如果微控制器具备输入输出内存保护单元（I/O
Memory Protection
Unit，IOMPU），也可以藉此限制DMA的访问范围，此时不需要独立的受信任保护域。

对于定时器等会自主发起高频中断的硬件，除了将其交给受信任保护域外，还可以在其中断向量中对中断数进行节流（Throttle），一旦发现中断频次过多则将其关闭，直到用户再次打开中断为止。当然，如果微控制器具备硬件节流能力，也可以藉此限制中断的频次。

\section{中断向量到保护域的划分}\label{sec:6-3}

RVM的中断分发链条从内核开始，一路向下，直到达到虚拟机内部。在这个链条上，越靠近内核的部分实时性越好、支持的中断频率越高，但安全性越低，对中断向量的可靠性要求越高。因此，对于某个特定频率的中断向量，需要选择合适的位置来放置它们。本小节将中断按照发生频率分为低频率、中等频率、高频率和超高频率四个类别，如下表所示\footnote{部分中断源虽然频率不高，但要求延迟很低，此时应当按延迟来计算其频率；比如某中断虽然每毫秒发生一次，但要求发生后在1微秒内作出响应，则其中断频率应该按照1MHz处理}。

\begin{xltabular}{\textwidth}{@{}>{\hsize=1.0000\hsize\raggedright\arraybackslash}X>{\hsize=1.0000\hsize\raggedright\arraybackslash}X>{\hsize=1.0000\hsize\raggedright\arraybackslash}X>{\hsize=1.0000\hsize\raggedright\arraybackslash}X>{\hsize=1.0000\hsize\raggedright\arraybackslash}X@{}}
\caption{中断按发生频率分类}\label{tab:6-2}\\
\toprule
类别 & 发生频率 & 间隔时间 & 代表性中断源 & 推荐划分方式\tabularnewline
\endfirsthead
\toprule
类别 & 发生频率 & 间隔时间 & 代表性中断源 & 推荐划分方式\tabularnewline
\midrule
\endhead
\bottomrule
\endlastfoot
低频率 & 1kHz以内 & 1毫秒以上；100k 周期以上 & 触摸屏中断 &
虚拟机中断/线程\tabularnewline
中等频率 & 1-10kHz & 1毫秒-100微秒；10k 周期以上 & IIC通信中断 &
原生进程内\tabularnewline
高频率 & 10-100kHz & 10微秒-100微秒；1k 周期以上 & SPI通信中断 &
内核可感知中断\tabularnewline
超高频率 & 1MHz以上 & 1微秒一下；100周期以内 &
超高波特率\footnote{1M及以上}串口中断 & 内核透明中断\tabularnewline
\end{xltabular}

\subsection{低频率中断向量}\label{subsec:6-3-1}

此种中断向量对性能要求很低，因此可以留在虚拟机内。对于移植过来的代码，其大体结构无需做任何更改，只要将其中断确认处理移动到内核内即可，而中断的大体则可留在虚拟机内；当这样处理时，中断延迟一般在几千个机器周期，也即十几微秒左右。这些中断一般是由人机交互设备或带缓冲区的通信设备产生的，无需很激进的响应时间也能达到不错的效果。

在安全性方面，此种中断向量的执行受虚拟机监视器的完全管辖，无论其发生任何错误都不会对虚拟机之外造成任何影响。

\subsection{中等频率中断向量}\label{subsec:6-3-2}

此种中断向量对性能有一定要求，将它们留在虚拟机内可能会导致中断丢失。因此，将它们移动到原生进程中是必要的；当这样处理时，中断延迟一般在几百个机器周期，也即几微秒左右。具体地，可以在原生进程中声明该向量，工程生成器会在内核中生成该向量对应的处理桩，并在原生进程中生成该向量对应的信号端点。在处理桩中需要进行的中断确认，并在桩返回时指示RME\_RVM\_VCT\_SIG\_SELF以仅向其专用端点发送信号。原生进程中的端点则会接收到该信号并解除中断向量的阻塞。

在安全性方面，此种中断向量的执行仍受内核提供的保护域限制，但不再受到虚拟机监视器的完全管辖，仅具备秘密性而不具备完整性和可用性，因此需要审查其代码以确保不会引入安全漏洞。

\subsection{高频率中断向量}\label{subsec:6-3-3}

此种中断向量对性能提出了相当高的要求，将它们留在任何保护域内都不可能，因为到保护域的中断延迟最少也需要数个微秒。因此，只能将这些中断向量的最关键部分放在内核中，并通过向原生进程发送信号或向虚拟机发送事件来推迟不关键部分的处理；当这样处理时，中断延迟一般在几十个机器周期，也即几百纳秒左右，但整个中断向量可能运行几百个机器周期以完成向信号端点的发送；如果在该中断发生时其他中断向量正在执行，则该中断必须等待上一中断向量执行完毕，造成几百个机器周期的延迟。

在安全性方面，此种中断向量不受任何限制。如果代码不经严格审查，会导致严重的功能和信息安全问题。

\subsection{超高频率中断向量}\label{subsec:6-3-4}

此种中断向量往往对响应时间和发生频率都提出了极端要求，以至于从中进行任何信号端点发送都不可能，否则一来延迟超标，二来CPU占用过高，甚至会使用户态得不到执行，失去了传递信号给用户态的意义。有些时候，这些中断还必须能够抢占系统中包括内核执行在内的任何执行，以确保固定的延迟时间。此时，只能这些中断向量做成透明中断，并使其抢占优先级高于内核\footnote{透明中断的抢占优先级是}：它们在进入时仅仅压栈那些调用约定中由调用者保存的寄存器，然后即开始执行C代码，完成后不通知内核就立刻弹栈退出；当这样处理时，中断延迟一般在几十个机器周期，也即几百纳秒左右，且整个中断向量的运行时间也可以做到少于一百个机器周期，实现裸机级别的响应能力。

为了使超高频率中断在发生时能够设法通知到用户态，可以另设一个可感知中断向量，在本中断向量退出时通过软件中断\footnote{如ARMv7-M的STIR寄存器触发方式}触发它，并让它替代自己通知内核。当然，考虑到中断的超高频率，触发可感知中断的频率不能太高，否则会使可感知中断消耗全部的CPU；此时可以对中断实行累计计次或施加速率限制，使得多次超高频率中断才对应于一次可感知中断触发。

在安全性方面，此种中断向量不受任何限制。如果代码不经严格审查，会导致严重的功能和信息安全问题。

\section{保护域间通信方法的选择}\label{sec:6-4}

RVM在保护域之间提供三种通信方法，分别是迁移调用、信号端点和事件源。其中，迁移调用和信号端点只能在原生进程之间使用，事件源则适用于原生进程与虚拟机以及虚拟机与虚拟机之间的通信。本小节将分情况介绍推荐的通信方法。

\subsection{耦合紧密的原生进程}\label{subsec:6-4-1}

对于耦合紧密、性能敏感的的原生进程间通信，推荐使用迁移调用，它可以使一个进程中的线程迁移到另一个进程中来运行，且在进入（RVM\_Inv\_Act）和返回（RVM\_Inv\_Ret）时都可通过寄存器组携带少量的数据，无需通过共享内存传送。它速度最快，但可能会造成目标进程内的并发访问：老线程迁移进入目标进程后对数据结构的访问可能会被调度器选择的信息新线程打断，此时新线程也可以迁移进入目标进程，并在老线程的访问结束之前发起同步访问。

因此，如果要使用迁移调用，则必须自行处理可能存在的并发性：加锁，或使用合理的无锁数据结构。

\subsection{耦合松散的原生进程}\label{subsec:6-4-2}

对于耦合松散、性能不敏感的的原生进程间通信，推荐使用信号端点和共享内存。信号端点自身具备序列化能力，即便发送端（RVM\_Sig\_Snd）有多个线程，接收端（RVM\_Sig\_Rcv）仍可以是单个线程，这能极大地简化通信编程。

值得注意的是，信号端点仅具备接收端单线程阻塞能力，但不具备发送端阻塞能力，也不具备接收端多线程阻塞能力。所幸，这些不具备的能力可以使用原生信号端点来实现：如果需要发送端阻塞，则可以给每个发送线程准备一个接收端点，在接收端完成服务、返回结果后向相应的发送端端点返回信号，令发送端解除阻塞；如果需要接收端多线程阻塞能力，则可使用多个接收端端点，而发送端可以在发送前判断接收端线程是否忙来决定唤醒何接收端线程。当然，这也可以通过巧妙地安排服务器和客户端的优先级实现：如果服务器的优先级比客户端高\footnote{且它们都位于同一个CPU上，这对于RVM而言是天然满足的}，则一旦服务端被唤醒，客户端就自然立刻阻塞，上述两个问题即迎刃而解，而且客户端可以用发送的信号数目来指示自身的身份。

\subsection{原生进程和虚拟机}\label{subsec:6-4-3}

原生进程到虚拟机和虚拟机到原生进程的通信必须通过信号端点实现。虚拟机到原生进程的发送和原生进程之间的发送别无二致：虚拟机中存在一个发送端点，而原生进程中存在一个接收端点，一旦虚拟机中对该发送端点发起发送，原生进程中接收端点上的线程即解除阻塞。然而，原生进程到虚拟机的发送则不同，接收端虚拟机必须在启动代码中将自己的虚拟中断注册到事件源（RVM\_Hyp\_Vct\_Evt），然后锁定虚拟机的权限设置（RVM\_Hyp\_Vct\_Lck）以防后续篡改\footnote{由于虚拟机每次重启都会强制运行一段用户定制的可信启动代码来进行授权和锁定，因此自授权是安全的}，最后再将相应的中断向量入口绑定到该虚拟中断（RVM\_Virt\_Vct\_Reg）\footnote{内部的中断向量入口}；原生进程则可以向任意事件源进行发送（RVM\_Prc\_Evt\_Snd）\footnote{这不会破坏虚拟机的可用性，因为所有的原生进程本就比虚拟机的优先级高}。一旦发送端的原生进程发送信号，接收端虚拟机的虚拟中断将被触发。

\subsection{虚拟机和虚拟机}\label{subsec:6-4-4}

虚拟机和虚拟机之间的通信也通过事件源进行。发送端虚拟机可以先自行授予向某个事件源发送的权限（RVM\_Hyp\_Evt\_Add），然后再锁定虚拟机的权限设置。接收端则在启动代码中将自己的虚拟中断注册到事件源，然后同样锁定虚拟机的权限设置，最后再将相应的中断向量入口绑定到该虚拟中断。

\section{生成前设置}\label{sec:6-5}

在生成工程之前，必须向工程生成器提供足够的有效信息。有关工程生成器的使用，请参见\protect\hyperref[sec:5-1]{第5章}
；本小节仅介绍填充各项设置时的注意事项。

\subsection{总体设置}\label{subsec:6-5-1}

此小节将介绍工程总体设置的推荐方法。

\subsubsection{设置基础选项}\label{subsubsec:6-5-1-1}

``基础选项''包括``工程名称''、``启用断言''和``启用调试打印''三项，其中``工程名称''可按需求决定。``启用断言''会启用内核内部的加严边界条件检查，``启用调试打印''则会在调试打印口上输出信息。调试打印口一般是一个串口，该串口会在芯片的RME启动头文件中指定。

\subsubsection{设置工作空间}\label{subsubsec:6-5-1-2}

``工作空间''包括``工作空间构建系统''和``工作空间输出覆盖''。这两项控制是否生成一个统一的工作空间，以及该工作空间是否需要被覆盖。``工作空间''的位置则在生成器的命令行中指定，或者在其图形界面的设置中指定。

很多构建系统，如Keil和Makefile等，都可以生成一个顶层工作空间，一旦在顶层工作空间中下令构建，各个子工作空间的构建便同步开始。要使用这个功能，需要在``工作空间构建系统''中选择一个构建系统，而且后续配置文件中内核、虚拟机监视器、原生进程和虚拟机的构建系统都必须与这个相同，否则这个功能处只能选择``None''，也即不生成统一的高层次工作空间构建系统。

``工作空间输出覆盖''则是指，如果重新生成整个工程，是否要用新生成的工作空间文件来覆盖老版本的文件。如果老版本文件上已经有了一些希望保留的修改，用户可以选择不覆盖，事实上这也是图形界面配置器提供的默认选项。``输出覆盖''选项也适用于本章节的某些其他选项，在那些选项处我们将指出它们也有``输出覆盖''选项，而不再对该选项的意图加以解释。

\subsubsection{设置芯片信息}\label{subsubsec:6-5-1-3}

``芯片信息''包括芯片的``架构''、``类别''和``名称''，以及具体芯片的配置项。以STM32F767IGT6为例，其架构为ARMv7-M，因此架构应选择``A7M''；类型为``STM32F767''，具体的芯片全名则是``STM32F767IGT6''。要将``类别''和``名称''分开成两个设置项的原因是，某些芯片在硬件配置上很相似，区别仅仅在于封装不同\footnote{封装差异一般体现在尾缀上}，没有必要将它们分开到不同的类别中，因为每一个类别都会对应一个独立的芯片描述文件。芯片信息

在选定芯片后，还需要决定需要启用的协处理器，以及一些其他参数。``协处理器''一般指浮点处理器，如STM32F767的Cortex-M7内核就有FPV4\_SP、FPV5\_SP和FPV5\_DP三个协处理器指令集，选中的协处理器将在工程中启用；没有在这里选中的协处理器不允许在任何进程或虚拟机中使用。

``其他参数''则无一定之规，一般至少会包括定时器超时值设置。``定时器超时值''指RME内核系统嘀嗒定时器的计时周期数，一旦经过这个数字的定时器周期就会产生一个时钟中断。一般而言，推荐将该值设置为长于每毫秒一次中断，这样中断的频率不会过高而消耗大量系统性能。以STM32F767为例，其主频为216MHz，因此定时器超时值应当至少设为216000。定时器周期的最大值可以一直放宽到系统需要的最小时间粒度：假设系统中所有的延时时间都是以10毫秒为单位的，则定时器超时值可以设为10毫秒以最小化系统开销。

\subsubsection{设置片外内存}\label{subsubsec:6-5-1-4}

``片外内存''指微控制器的外扩存储器。部分微控制器具有片外存储器接口，这些片外存储器可以像片内存储器那样被访存指令直接访问。由于片外存储器并不在芯片的描述文件中，因此如果需要让工程生成器知晓它们的存在，就必须在此单独予以声明。

部分芯片没有片上ROM/RAM，此时所有的ROM/RAM都是片外内存，需要在创建工程时被单独添加到该位置，以便后续配置中使用。

\subsubsection{设置共享内存}\label{subsubsec:6-5-1-5}

``共享内存''指多个保护域都有权限访问的设备存储器。它们必须在共享内存列表中声明其基地址和长度，以及可能的最大权限，然后才能在各保护域的共享内存列表中被引用，以使各有关保护域有能力访问该段共享内存。在引用时的权限必须是最大权限的子集；这意味着，即便是同一段共享内存，不同的保护域对其的访问权限也可以不同。

各共享内存段的名称必须指定以便后续引用；除了设备段存储器基地址必须人工指定、其对齐方式必须为``Auto''之外，其余类别存储器的基地址和对齐方式都可以由工程生成器自行决定。

\subsubsection{设置页表方式}\label{subsubsec:6-5-1-6}

如果内存权限在整个过程中并不修改\footnote{绝大多数使用工程生成器的场合都是如此}，而且不使用动态映射页面\footnote{每个保护域声明的内存}，则可以打开``启用固定式裸页表''开关，在内核中除去页表逻辑；此时，页表将被ROM中的固定MPU数据表格代替，可以节约页表数据结构和管理代码所占的空间。

\subsection{内核设置}\label{subsec:6-5-2}

此小节将介绍RME内核的推荐设置方法。

\subsubsection{添加存储器}\label{subsubsec:6-5-2-1}

内核需要一段代码段存储器和一段数据段存储器用来存放自己。代码段起始地址推荐放在片上ROM地址空间的开头，那里一般还会存放中断向量表。代码段的大小只要以能放下内核并通过链接为限。数据段起始地址同样推荐放在片上RAM地址空间的开头\footnote{如果芯片具备紧耦合内存（Tightly-Coupled
  Memory，TCM），推荐放在那里以最大化性能}，大小也以能通过编译为限。数据段中包含的除了内核本身的固定数据结构之外，还包含有内核对象、内核栈等，但这些结构已经被生成的链接器脚本考虑在内，只要链接能通过，RAM的分配量就是合适的。

为了简化用户的心智模型和提高系统的可靠性，在所有的官方标准移植中，RME在执行时总是关闭内存保护单元从而使用默认的内存映射属性\footnote{也即经典裸机程序使用的映射属性；比如，在ARMv7-M中，0x2000000开头的片上内存默认是RWBXCS的}。用户必须保证代码段和数据段在默认映射下可用，且具备合适的权限和缓存属性；这在所有现时微控制器的片内内存上都是成立的\footnote{为了方便裸机编程，厂家通常这样设计默认的内存映射属性}。

\subsubsection{内核参数设置}\label{subsubsec:6-5-2-2}

``内核栈长度''是内核所使用的栈的大小。该值设为1kB一般就足够了；如果自定义中断上半部占用了过大的栈，则可以改为4kB。

``额外内核对象用内存长度''一般配置为0，除非需要手动创建内核对象；这只在全定制时才需要。

``内核内存分配粒度''是内核对象创建时的粒度的2的次方。粒度越大，分配就越粗糙，但是记录某块内存分配与否的内核内存登记表的所占空间就越小。常见的内核对象大都在16字节以内。

对于小型系统，推荐设为3，此时将以8字节的粒度分配内核内存；对于较复杂、耗用内核内存较多的场合可以设为4，此时将以16字节的粒度分配内核内存以减少内核内存登记表的大小；如果系统极为复杂，设为5甚至是6也是可以的，此时粒度为32或64。

\subsubsection{工具链设置}\label{subsubsec:6-5-2-3}

``工具链设置''中可以设置``工具链''、``构建系统''以及``优化等级''。``工具链''和``构建系统''有联动关系，并非所有的工具链和构建系统组合都可以选择。一旦选择了某个工具链，只有与之相兼容的构建系统才会出现在构建系统选项中，且工程一旦生成则无法后需更改。``优化等级''则可自由选择合适的级别。

\subsubsection{文件路径设置}\label{subsubsec:6-5-2-4}

``文件路径设置''中设置各类文件的相对路径。该相对路径是对子工程的文件夹而言的，它们会在工程创建时一并被创建。如果没有特殊要求，这些路径使用工程生成器图形界面中给出的默认值即可，无需更改。

\subsection{虚拟机监视器设置}\label{subsec:6-5-3}

此小节将介绍RVM虚拟机监视器的推荐设置方法。

\subsubsection{添加存储器}\label{subsubsec:6-5-3-1}

虚拟机监视器也需要一段代码段存储器和一段数据段存储器用来存放自己。不过，与内核不同，其代码段和数据段的起始地址将直接放在内核之后，因此可以配置的仅仅包括其代码段和数据段的大小；该大小也以能通过编译为限，只要链接能通过，存储器的分配量就是合适的。

为了适应各种微控制器架构，虚拟机监视器本身总是将整个地址空间都映射为可缓存和可缓冲属性，包括那些潜在的设备区域。这是没有问题的，因为虚拟机监视器不会访问任何设备\footnote{如果用户主动在虚拟机监视器的钩子中插入访问设备的代码，则必须注意内存访问属性；为了降低用户的心智负担，不建议在虚拟机监视器的钩子（如物理中断源和事件源的过滤器等）中访问设备}。然而，虚拟机监视器需要访问一段存有中断源和事件源标志位的内核内存，该内核内存在内核中则使用背景区域属性。因此，用户必须保证内核内存也是在背景区域属性下也是可缓存和可缓冲的。

\subsubsection{虚拟化参数设置}\label{subsubsec:6-5-3-2}

与内核不同，虚拟机监视器需要配置的栈长度有三个，分别是``启动守护线程栈长度''、``安全守护线程栈长度''和``虚拟机守护线程栈长度''，这是因为虚拟机监视器由三个线程组成所致。``启动守护线程''和``安全守护线程''的职责较为简单，256B的栈一般是足够的；虚拟机守护线程则至少需要512B的栈。在内存空间充裕的场合，栈预算可以翻番。

``额外权能表槽位数''一般配置为0，除非需要手动创建内核对象；这只在全定制时才需要。

``启用空闲休眠''指示虚拟机监视器的``启动守护线程''在空闲时将处理器置于休眠模式。该项一般选择Yes；如果不希望进入低功耗休眠状态以尽量降低中断延迟，可以选择``No''。

``虚拟化优先级数''一般配置为一个机器字长，因为RVM支持的处理器至少是32位，而32个优先级对一般的虚拟化系统都足够了。如果不使用虚拟化功能，该参数可以配置为0。

``虚拟机事件源数''与``虚拟机向量映射数''可以按需配置，使用多少就配置多少。事件源的数量必须是处理器字长的整数倍且不能超过1024，映射数则无此限制。如果不使用虚拟化功能，两者都可以配置为0。

\subsubsection{工具链设置}\label{subsubsec:6-5-3-3}

该项设置与内核的设置无区别，请参见内核设置的相关章节。

\subsubsection{文件路径设置}\label{subsubsec:6-5-3-4}

该项设置与内核的设置无太大区别，推荐保持默认选项不变。

\subsection{原生进程设置}\label{subsec:6-5-4}

此小节将介绍原生进程的推荐设置方法。

\subsubsection{添加私有存储器}\label{subsubsec:6-5-4-1}

原生进程的私有存储器包括至少一个代码段存储器、至少一个数据段存储器和任意数量的设备段存储器。在编译时，其代码段和数据段总是分别放在列表中的第一个代码段存储器和数据段存储器上。除了设备段存储器基地址必须人工指定、其对齐方式必须为``Auto''之外，其余类别存储器的基地址和对齐方式都可以由工程生成器自行决定。如果不指定名称，则该段存储器为匿名存储器，其基地址和长度的宏将不可用。

\subsubsection{添加共享存储器}\label{subsubsec:6-5-4-2}

原生进程可以包含也可以不包含共享存储器。要引用共享存储器，只要填写它们的名称，以及它们在此进程中的权限即可。共享存储器的基地址、长度、对齐方式和最大权限的设置见\protect\hyperref[subsubsec:6-5-1-5]{6.5.1.5}
。

\subsection{虚拟机设置}\label{subsec:6-5-5}

此小节将介绍虚拟机的推荐设置方法。

\subsubsection{添加私有存储器}\label{subsubsec:6-5-5-1}

和原生进程类似，虚拟机的私有存储器也包括至少一个代码段存储器、至少一个数据段存储器和任意数量的设备段存储器。在编译时，其代码段和数据段同样总是分别放在列表中的第一个代码段存储器和数据段存储器上。然而，由于虚拟机内部Usr线程的寄存器组内存是与内核共享的，因此虚拟机的第一个数据段的访问属性必须与默认的内存映射属性一致。

\subsubsection{添加共享存储器}\label{subsubsec:6-5-5-2}

共享存储器的添加方法与原生进程相同。
